% figures/w03-why-T.svg — 배분행렬 T 가 왜 [1 1; b -b] 인가. 열을 하나씩 세운다.
%
%   bash _tools/tikz2svg.sh figures/src/w03-why-T.tex figures/w03-why-T.svg
%
% 사용자 지시 (2026-10-13): "왜 T 행렬이 저렇게 되는지도 선박 그림과 함께
% (다른 곳 보면 있을 거야..) 다시 찾아서 자세히 그려줘봐"
%
% "다른 곳" 은 부록 A1 이다. `figures/a1-column-rule.svg` 가 **일반 규칙**을 유도한다 —
% (x,y) 에서 e 방향으로 미는 추진기 하나가 B 의 열 하나를 만든다:
%
%     b_i = [ e_x ; e_y ; x_i e_y - y_i e_x ]^T        (A1 2절)
%
% 이 그림은 그 규칙을 **Otter 의 두 추진기에 실제로 적용**해서 T 를 세운다. A1 의 그림이
% "규칙이 왜 맞는가" 라면, 이 그림은 "그 규칙이 이 배에서 무엇을 내놓는가" 다.
% 3 주차 학생은 A1 을 아직 읽지 않았으므로 규칙도 함께 적되, 유도는 A1 로 넘긴다.
%
% 수치는 MATLAB 으로 확인했다 — otter_B 와 0 차이 (2026-10-13).
%
% 배치 주의 (렌더해서 고친 것):
%   선체 패널은 **글자를 선체 밖으로** 빼야 한다. 폰툰이 y = -1.7 .. 1.25 를 채우므로
%   위치 벡터의 좌표·지레·회전 설명이 전부 그 위에 겹쳤다. 좌표는 폰툰 아래,
%   지레는 더 아래, 회전 설명은 선체 오른쪽 바깥에 둔다.
\documentclass[border=6pt]{standalone}
% gnc-style.tex — 이 볼트의 TikZ 그림이 공유하는 스타일.
%
% 저널 그림 규약을 스타일 하나로 굳혀 둔다. 그림마다 선 굵기나 화살촉을
% 다시 정하지 않는다 — 그것이 그림이 제각각으로 보이는 원인이다.
%
% 쓰는 법:  % gnc-style.tex — 이 볼트의 TikZ 그림이 공유하는 스타일.
%
% 저널 그림 규약을 스타일 하나로 굳혀 둔다. 그림마다 선 굵기나 화살촉을
% 다시 정하지 않는다 — 그것이 그림이 제각각으로 보이는 원인이다.
%
% 쓰는 법:  % gnc-style.tex — 이 볼트의 TikZ 그림이 공유하는 스타일.
%
% 저널 그림 규약을 스타일 하나로 굳혀 둔다. 그림마다 선 굵기나 화살촉을
% 다시 정하지 않는다 — 그것이 그림이 제각각으로 보이는 원인이다.
%
% 쓰는 법:  \input{gnc-style}  (figures/src/ 안에서)

\usepackage{tikz}
\usepackage{amsmath}
\usepackage{xcolor}
\usetikzlibrary{arrows.meta, positioning, calc, fit, backgrounds, decorations.pathmorphing}

% ---- 색 --------------------------------------------------------------------
% 색은 강조 하나에만 쓴다. 나머지는 검정.
\definecolor{ink}{HTML}{1A1A1A}
\definecolor{accent}{HTML}{7C3AED}
\definecolor{warn}{HTML}{B91C1C}
\definecolor{soft}{HTML}{666666}

% 기하 그림(좌표계·유도)에서만 쓰는 넷. 블록선도는 위의 흑백 + accent 로 충분하다.
\definecolor{pathc}{HTML}{0D9488}   % 계획 경로 · 항적
\definecolor{errc}{HTML}{C2410C}    % 오차 (cross-track 등)
\definecolor{alng}{HTML}{B45309}    % along-track
\definecolor{flow}{HTML}{0369A1}    % 조류 · 유속

% ---- 선과 화살촉 -----------------------------------------------------------
\tikzset{
  % 신호선. 화살촉은 작고 뾰족한 것 하나뿐이다.
  sig/.style   = {ink, line width=0.5pt, -{Stealth[length=4pt, width=3pt]}},
  wire/.style  = {ink, line width=0.5pt},
  % 강조 경로 — 파선으로 구분한다. 흑백 인쇄에서도 살아야 하기 때문이다.
  asig/.style  = {accent, line width=0.5pt, densely dashed,
                  -{Stealth[length=4pt, width=3pt]}},
  awire/.style = {accent, line width=0.5pt, densely dashed},
  % 블록. 흰 바탕, 직각, 검은 테두리.
  blk/.style   = {draw=ink, line width=0.55pt, fill=white,
                  minimum width=17mm, minimum height=8mm, inner sep=2pt, font=\small},
  % 합산점
  sum/.style   = {draw=ink, line width=0.55pt, fill=white, circle,
                  inner sep=0pt, minimum size=4.6mm, font=\scriptsize},
  asum/.style  = {draw=accent, line width=0.55pt, fill=white, circle,
                  inner sep=0pt, minimum size=4.6mm, font=\scriptsize},
  % 분기점
  dot/.style   = {ink, fill=ink, circle, inner sep=0pt, minimum size=1.5mm},
  % 신호 이름 — 선 위에 작게
  sname/.style = {font=\small, inner sep=1.5pt},
  % 그림 안의 짧은 설명. 그림 안의 글은 최소로 둔다
  note/.style  = {font=\scriptsize, soft, inner sep=1.5pt},
  anote/.style = {font=\scriptsize, accent, inner sep=1.5pt},
  % 축
  axis/.style  = {ink, line width=0.5pt},
  tick/.style  = {font=\scriptsize, soft},
}

% 캡션 한 줄 — 그림 아래 가는 선과 함께
\newcommand{\gnccaption}[2]{%
  \node[anchor=north west, text width=#1, align=left, font=\scriptsize, ink]
        at (current bounding box.south west) {\vspace{2pt}#2};}
  (figures/src/ 안에서)

\usepackage{tikz}
\usepackage{amsmath}
\usepackage{xcolor}
\usetikzlibrary{arrows.meta, positioning, calc, fit, backgrounds, decorations.pathmorphing}

% ---- 색 --------------------------------------------------------------------
% 색은 강조 하나에만 쓴다. 나머지는 검정.
\definecolor{ink}{HTML}{1A1A1A}
\definecolor{accent}{HTML}{7C3AED}
\definecolor{warn}{HTML}{B91C1C}
\definecolor{soft}{HTML}{666666}

% 기하 그림(좌표계·유도)에서만 쓰는 넷. 블록선도는 위의 흑백 + accent 로 충분하다.
\definecolor{pathc}{HTML}{0D9488}   % 계획 경로 · 항적
\definecolor{errc}{HTML}{C2410C}    % 오차 (cross-track 등)
\definecolor{alng}{HTML}{B45309}    % along-track
\definecolor{flow}{HTML}{0369A1}    % 조류 · 유속

% ---- 선과 화살촉 -----------------------------------------------------------
\tikzset{
  % 신호선. 화살촉은 작고 뾰족한 것 하나뿐이다.
  sig/.style   = {ink, line width=0.5pt, -{Stealth[length=4pt, width=3pt]}},
  wire/.style  = {ink, line width=0.5pt},
  % 강조 경로 — 파선으로 구분한다. 흑백 인쇄에서도 살아야 하기 때문이다.
  asig/.style  = {accent, line width=0.5pt, densely dashed,
                  -{Stealth[length=4pt, width=3pt]}},
  awire/.style = {accent, line width=0.5pt, densely dashed},
  % 블록. 흰 바탕, 직각, 검은 테두리.
  blk/.style   = {draw=ink, line width=0.55pt, fill=white,
                  minimum width=17mm, minimum height=8mm, inner sep=2pt, font=\small},
  % 합산점
  sum/.style   = {draw=ink, line width=0.55pt, fill=white, circle,
                  inner sep=0pt, minimum size=4.6mm, font=\scriptsize},
  asum/.style  = {draw=accent, line width=0.55pt, fill=white, circle,
                  inner sep=0pt, minimum size=4.6mm, font=\scriptsize},
  % 분기점
  dot/.style   = {ink, fill=ink, circle, inner sep=0pt, minimum size=1.5mm},
  % 신호 이름 — 선 위에 작게
  sname/.style = {font=\small, inner sep=1.5pt},
  % 그림 안의 짧은 설명. 그림 안의 글은 최소로 둔다
  note/.style  = {font=\scriptsize, soft, inner sep=1.5pt},
  anote/.style = {font=\scriptsize, accent, inner sep=1.5pt},
  % 축
  axis/.style  = {ink, line width=0.5pt},
  tick/.style  = {font=\scriptsize, soft},
}

% 캡션 한 줄 — 그림 아래 가는 선과 함께
\newcommand{\gnccaption}[2]{%
  \node[anchor=north west, text width=#1, align=left, font=\scriptsize, ink]
        at (current bounding box.south west) {\vspace{2pt}#2};}
  (figures/src/ 안에서)

\usepackage{tikz}
\usepackage{amsmath}
\usepackage{xcolor}
\usetikzlibrary{arrows.meta, positioning, calc, fit, backgrounds, decorations.pathmorphing}

% ---- 색 --------------------------------------------------------------------
% 색은 강조 하나에만 쓴다. 나머지는 검정.
\definecolor{ink}{HTML}{1A1A1A}
\definecolor{accent}{HTML}{7C3AED}
\definecolor{warn}{HTML}{B91C1C}
\definecolor{soft}{HTML}{666666}

% 기하 그림(좌표계·유도)에서만 쓰는 넷. 블록선도는 위의 흑백 + accent 로 충분하다.
\definecolor{pathc}{HTML}{0D9488}   % 계획 경로 · 항적
\definecolor{errc}{HTML}{C2410C}    % 오차 (cross-track 등)
\definecolor{alng}{HTML}{B45309}    % along-track
\definecolor{flow}{HTML}{0369A1}    % 조류 · 유속

% ---- 선과 화살촉 -----------------------------------------------------------
\tikzset{
  % 신호선. 화살촉은 작고 뾰족한 것 하나뿐이다.
  sig/.style   = {ink, line width=0.5pt, -{Stealth[length=4pt, width=3pt]}},
  wire/.style  = {ink, line width=0.5pt},
  % 강조 경로 — 파선으로 구분한다. 흑백 인쇄에서도 살아야 하기 때문이다.
  asig/.style  = {accent, line width=0.5pt, densely dashed,
                  -{Stealth[length=4pt, width=3pt]}},
  awire/.style = {accent, line width=0.5pt, densely dashed},
  % 블록. 흰 바탕, 직각, 검은 테두리.
  blk/.style   = {draw=ink, line width=0.55pt, fill=white,
                  minimum width=17mm, minimum height=8mm, inner sep=2pt, font=\small},
  % 합산점
  sum/.style   = {draw=ink, line width=0.55pt, fill=white, circle,
                  inner sep=0pt, minimum size=4.6mm, font=\scriptsize},
  asum/.style  = {draw=accent, line width=0.55pt, fill=white, circle,
                  inner sep=0pt, minimum size=4.6mm, font=\scriptsize},
  % 분기점
  dot/.style   = {ink, fill=ink, circle, inner sep=0pt, minimum size=1.5mm},
  % 신호 이름 — 선 위에 작게
  sname/.style = {font=\small, inner sep=1.5pt},
  % 그림 안의 짧은 설명. 그림 안의 글은 최소로 둔다
  note/.style  = {font=\scriptsize, soft, inner sep=1.5pt},
  anote/.style = {font=\scriptsize, accent, inner sep=1.5pt},
  % 축
  axis/.style  = {ink, line width=0.5pt},
  tick/.style  = {font=\scriptsize, soft},
}

% 캡션 한 줄 — 그림 아래 가는 선과 함께
\newcommand{\gnccaption}[2]{%
  \node[anchor=north west, text width=#1, align=left, font=\scriptsize, ink]
        at (current bounding box.south west) {\vspace{2pt}#2};}


\definecolor{knownc}{HTML}{1D4ED8}

\tikzset{
  hull/.style    = {ink, line width=0.9pt, fill=black!4},
  pont/.style    = {ink, line width=0.7pt, fill=pathc!12},
  evec/.style    = {errc, line width=1.5pt, -{Stealth[length=5.2pt,width=3.8pt]}},
  rvec/.style    = {knownc, line width=1.1pt, -{Stealth[length=4.6pt,width=3.2pt]}},
  ctr/.style     = {soft, line width=0.4pt, densely dashed},
  spin/.style    = {alng, line width=1.2pt, -{Stealth[length=4.6pt,width=3.4pt]}},
  ax/.style      = {ink, line width=0.7pt, -{Stealth[length=4.6pt,width=3.2pt]}},
  lever/.style   = {knownc, line width=0.8pt,
                    {Stealth[length=3.2pt,width=2.2pt]}-{Stealth[length=3.2pt,width=2.2pt]}},
  lbl/.style     = {font=\scriptsize},
  rule/.style    = {draw=accent, line width=1.0pt, fill=accent!5, inner sep=5pt,
                    rounded corners=2pt, align=center, font=\small},
  colL/.style    = {draw=errc, line width=0.9pt, fill=errc!5, inner sep=5pt,
                    rounded corners=2pt, align=left, font=\scriptsize},
  colR/.style    = {draw=flow, line width=0.9pt, fill=flow!5, inner sep=5pt,
                    rounded corners=2pt, align=left, font=\scriptsize},
  want/.style    = {draw=accent, line width=1.1pt, fill=white, inner sep=5pt,
                    rounded corners=2pt, align=center, font=\small},
  capt/.style    = {font=\scriptsize\bfseries, warn},
  stage/.style   = {font=\small\bfseries, ink},
}

\begin{document}
\begin{tikzpicture}[x=1cm, y=1cm]

% =====================================================================
% STEP 0 — 규칙 (부록 A1 에서 유도한 것)
% =====================================================================
\node[stage, anchor=west] at (-2.8,4.55) {STEP 0 \ \ the rule, derived in Appendix A1};
\node[rule, text width=8.2cm, anchor=north west] at (-2.8,4.15)
  {one propeller at $(x_i, y_i)$ pushing along the unit vector $\mathbf{e}_i$
   contributes \textbf{one column} of $\mathbf{B}$\\[5pt]
   $\mathbf{b}_i = \big[\;e_x \quad e_y \quad x_i e_y - y_i e_x \;\big]^{\!\top}$\\[4pt]
   {\scriptsize the third entry is the $z$ component of $\mathbf{r}_i \times \mathbf{e}_i$ ---
    force times perpendicular distance, with the sign already in it}};

% =====================================================================
% STEP 1 — 그 규칙을 이 배에 적용한다
% =====================================================================
\node[stage, anchor=west] at (-2.8,1.25) {STEP 1 \ \ apply it to this hull, one propeller at a time};

\begin{scope}[shift={(0.4,-3.1)}]
  \draw[ctr] (0,-2.9) -- (0,3.1);
  \draw[pont] (-1.75,-1.7) rectangle (-0.95,1.25);
  \draw[pont] ( 0.95,-1.7) rectangle ( 1.75,1.25);
  \draw[hull] (-1.75,0.35) rectangle (1.75,1.25);

  %  축 — 이름표는 화살표 끝 바깥에 / axes, named beyond their tips
  \draw[ax] (0,0) -- (0,2.95);   \node[lbl, anchor=south] at (0,3.02) {$x_b$ \ forward};
  \draw[ax] (0,0) -- (3.05,0);   \node[lbl, anchor=west]  at (3.12,0) {$y_b$ \ starboard};
  \fill[ink] (0,0) circle (1.3pt);
  \node[lbl, anchor=north east] at (-0.08,-0.10) {$o_b$};

  %  위치 벡터 r / the position vectors
  \draw[rvec] (0,0) -- (-1.35,0);
  \draw[rvec] (0,0) -- ( 1.35,0);
  \node[lbl, knownc, anchor=south] at (-0.70,0.05) {$\mathbf{r}_L$};
  \node[lbl, knownc, anchor=south] at ( 0.70,0.05) {$\mathbf{r}_R$};
  %  좌표는 폰툰 **아래** / the coordinates go below the pontoons
  \node[lbl, knownc, anchor=north] at (-1.35,-1.80) {$(0,\,-b)$};
  \node[lbl, knownc, anchor=north] at ( 1.35,-1.80) {$(0,\,+b)$};

  %  추진 방향 e — 둘 다 앞으로 / both unit vectors point forward
  \draw[evec] (-1.35,1.25) -- (-1.35,2.45);
  \draw[evec] ( 1.35,1.25) -- ( 1.35,2.45);
  \node[lbl, errc, anchor=south east] at (-1.20,2.50) {$\mathbf{e}_L$};
  \node[lbl, errc, anchor=south west] at ( 1.50,2.50) {$\mathbf{e}_R$};
  \node[lbl, errc, anchor=east, align=right] at (-1.95,1.95)
    {both propellers are\\ \textbf{fixed} and face\\ forward: $\mathbf{e} = (1,0)$};

  %  지레 — 폰툰보다 더 아래 / the lever, below the coordinates
  \draw[lever] (-1.35,-2.45) -- (0,-2.45);
  \node[lbl, knownc, anchor=north] at (-0.68,-2.52) {$b = 0.3950$ m};

  %  회전은 선체 **오른쪽 위**의 빈 곳에. 아래에 두면 가운데 단락과 겹친다 (렌더해서 고침)
  \draw[spin] (1.95,1.95) arc (62:14:2.3);
  \node[lbl, alng, anchor=west, align=left] at (2.25,1.55)
    {push \textbf{port} alone\\
     $\Rightarrow$ bow swings\\
     \ \ \ to starboard\\
     $\Rightarrow$ $N > 0$};

  \node[capt, anchor=north] at (0.4,-3.25) {seen from above};
\end{scope}

% =====================================================================
% 열 두 개를 세운다
% =====================================================================
\node[colL, text width=5.6cm, anchor=north west] at (5.6,1.00)
  {\textbf{\textcolor{errc}{port propeller}} \ --- \ $(x,y) = (0,\,-b)$\\[4pt]
   $e_x = 1$ \hfill all of its push goes into $X$\\[2pt]
   $e_y = 0$ \hfill none of it goes sideways\\[2pt]
   $x e_y - y e_x = 0\cdot 0 - (-b)\cdot 1 = \mathbf{+b}$\\[4pt]
   $\Rightarrow\ \mathbf{b}_L = \big[\,1 \quad 0 \quad +b\,\big]^{\!\top}$};

\node[colR, text width=5.6cm, anchor=north west] at (5.6,-1.75)
  {\textbf{\textcolor{flow}{starboard propeller}} \ --- \ $(x,y) = (0,\,+b)$\\[4pt]
   $e_x = 1$ \hfill the same push\\[2pt]
   $e_y = 0$ \hfill the same nothing sideways\\[2pt]
   $x e_y - y e_x = 0\cdot 0 - (+b)\cdot 1 = \mathbf{-b}$\\[4pt]
   $\Rightarrow\ \mathbf{b}_R = \big[\,1 \quad 0 \quad -b\,\big]^{\!\top}$};

\node[anchor=north west, align=left, font=\scriptsize, soft, text width=5.6cm] at (5.6,-4.50)
  {\textbf{only the sign of $y$ differs}, and it is the one that
   matters: the two sit on opposite sides, so their moments
   are opposite and their pushes are not};

% =====================================================================
% STEP 2 — 열을 나란히 세우면 B, 0 행을 지우면 T
% =====================================================================
\node[stage, anchor=west] at (12.0,1.25) {STEP 2 \ \ stack the columns};

\node[want, anchor=north west] at (12.0,1.00)
  {$\mathbf{B} = \big[\ \mathbf{b}_L \ \ \mathbf{b}_R\ \big]
    = \begin{bmatrix} 1 & 1 \\ 0 & 0 \\ b & -b \end{bmatrix}
    \begin{matrix} \scriptstyle\leftarrow X \\ \scriptstyle\leftarrow Y \\ \scriptstyle\leftarrow N \end{matrix}$};

\node[anchor=north west, align=left, font=\scriptsize, warn, text width=5.1cm] at (12.0,-0.90)
  {\textbf{the middle row is empty.} No combination of
   $F_L$ and $F_R$ makes $Y$, so this hull \emph{cannot be
   pushed sideways} --- which is why every later week
   turns the vessel in order to move it};

\node[want, anchor=north west] at (12.0,-2.60)
  {delete that row:\\[5pt]
   $\begin{bmatrix} X \\ N \end{bmatrix}
    = \underbrace{\begin{bmatrix} 1 & 1 \\ b & -b \end{bmatrix}}_{\mathbf{T}}
      \begin{bmatrix} F_L \\ F_R \end{bmatrix}$};

\node[anchor=north west, align=left, font=\scriptsize, text width=5.1cm] at (12.0,-4.50)
  {$2\times 2$ with $\det\mathbf{T} = -2b \neq 0$: two actuators,
   two controllable forces, and exactly one way to
   split them --- so the inverse of \S3-6 exists};

\node[capt, anchor=north, align=center] at (14.5,-5.85)
  {$\mathbf{T}$ is not chosen --- it is read off the hull};

\end{tikzpicture}
\end{document}
